\documentclass{../yalumba}

\title{Spectral Ecosystem Inference for\\Reference-Free Relatedness Detection:\\From Graph Collapse to Multi-Scale Signatures}
\author{Ajax Davis}
\date{April 2026}

\setabstract{%
We present the design, implementation, and experimental evaluation of
\emph{Spectral Ecology}---the first nonlinear, identity-free algorithm in the
yalumba symbiogenesis family.  Where prior algorithms (Module Persistence,
Coalition Transfer, Module Stability) compared genomes by measuring overlap of
shared tokens, Spectral Ecology asks whether two genomes instantiate the same
\emph{kind} of organization.  The method constructs dense module co-occurrence
graphs (600--1{,}400 edges from 295--449 modules per sample), computes their
normalized graph Laplacians via Jacobi eigendecomposition, extracts multi-scale
heat kernel signatures at 10 timescales, assigns spectral ecological roles
(core, satellite, bridge, scaffold, boundary), estimates eigenvalue density
profiles, eigenvector localization statistics, and diffusion distance
distributions---then compares these identity-free invariants across samples.
We document the complete experimental trajectory: v1 failed catastrophically
because sequential adjacency produced graphs that were $>$99\% disconnected
(2--5 edges, spectral gap $= 0$); v2 switched to read-level co-occurrence,
producing 250$\times$ denser graphs and flipping separation from $-0.74\%$ to
$+1.01\%$; five subsequent experiments iterated on connectivity guarantees,
component filtering, eigenvalue filtering, node-level HKS distribution
comparison, and mean-normalization, pushing separation to $+1.89\%$.  On the
CEPH~1463 six-member pedigree benchmark, the method achieves $+1.55\%$ mean
separation with three of four parent--child pairs ranking consecutively at
positions 3--5.  Two persistent failure modes are diagnosed in detail:
(1)~the spouse pair (Father~$\leftrightarrow$~Mother) consistently ranks first
because their ecosystems have nearly identical module counts, edge densities,
and spectral shapes---a coverage confound that purely within-sample spectral
methods cannot resolve; (2)~NA12891 (Mat.~GF) is a structural outlier with a
unique role distribution (356 bridges vs.\ 0 in all other samples), dragging
the Mat.GF~$\to$~Mother parent--child pair to rank~8.  We argue that the
spectral approach has reached its ceiling for within-sample comparison
($\sim$1.5--1.9\% separation) and that further progress requires cross-sample
structural signal---multi-layer ecosystem graphs, module embedding similarity,
or optimal transport between spectral distributions.  The \texttt{@yalumba/ecology}
package implementing all methods is released as part of the yalumba monorepo:
Jacobi eigendecomposition, multi-scale heat kernel, spectral role assignment,
eigenvalue density estimation, eigenvector localization, null model
deconfounding, and module feature embeddings---all from scratch in TypeScript
with no external dependencies.}

\begin{document}
\maketitle

\tableofcontents
\newpage


% ════════════════════════════════════════════════════════════════════
\section{Introduction}
\label{sec:intro}

\subsection{The Identity Trap}

Every prior symbiogenesis algorithm fell into the same trap: they compared
genomes by asking ``do these two samples share the same tokens?''  Module
Persistence asked whether two samples share the same modules.  Coalition
Transfer asked whether they share the same coalitions.  Module Stability
asked whether the same motifs are stable under perturbation.

\ymarginnote{Token-matching fails because related individuals produce
related-but-not-identical structures.}

All three hit the same wall.  Related individuals who share a haplotype block
produce k-mer modules from that block, but heterozygosity at even one position
changes the k-mer composition, producing \emph{different} module identities.
The modules are biologically related but computationally distinct.  Exact
matching misses them entirely.

The weight-tuning ceiling was proven by Module Stability~v4, which
exhaustively explored six scoring configurations and proved that the coverage
confound---samples with similar sequencing depth produce similar module
structures regardless of kinship---is structural, not weight-tunable.

\subsection{The Spectral Alternative}

Spectral Ecology abandons token identity entirely.  Instead of asking ``do
these samples share tokens?'' it asks:

\begin{pullquote}
Do these two genomes instantiate the same \emph{kind} of genomic organization?
\end{pullquote}

The answer is computed from spectral geometry: eigenvalues and eigenvectors of
the graph Laplacian, heat kernel diffusion at multiple timescales, ecological
role distributions, and eigenvalue density profiles.  None of these depend on
which specific modules exist---only on how they are organized.

\ymarginnote{Spectral methods measure structural similarity without node
correspondence.}

Two graphs with identical spectra have identical topology regardless of node
labeling~\cite{fang2012selection}.  Two graphs with similar heat kernel traces
have similar multi-scale connectivity.  Two graphs with similar role
distributions have similar ecological architecture.  This is exactly the
property needed: related individuals may produce different modules, but the
\emph{organization} of those modules---the interaction pattern, the hub
structure, the clustering---should be inherited.

\subsection{Contributions}

This paper documents:

\begin{enumerate}
  \item The complete \texttt{@yalumba/ecology} package: graph Laplacian
    computation, Jacobi eigendecomposition, multi-scale heat kernel signatures,
    spectral ecological role assignment, eigenvalue density estimation,
    eigenvector localization, null model deconfounding, and module feature
    embeddings---all from scratch in TypeScript.
  \item The experimental trajectory from v1 (catastrophic graph collapse) to v2
    (dense co-occurrence graphs) through five iterations that pushed separation
    from $-0.74\%$ to $+1.89\%$.
  \item A detailed diagnosis of the two persistent failure modes: the spouse
    confound and the NA12891 outlier.
  \item An analysis of why purely within-sample spectral methods plateau at
    $\sim$1.5\% separation and what architectural changes are needed to break
    through.
\end{enumerate}

\subsection{Context: The Symbiogenesis Program}

Spectral Ecology is the fourth algorithm family in the symbiogenesis
program~\cite{margulis1998}.  The program's thesis:

\begin{definition}
Human inheritance can be modeled as the transmission of resilient, interacting
genomic ecosystems whose structure can be inferred from raw reads without
reference coordinates.  Close kinship is only one observable of this deeper
system.
\end{definition}

The program operates on two tracks:

\begin{itemize}
  \item \textbf{Track~A (practical baseline):} Rare-Run~P90---contiguous runs
    of rare k-mers.  Separation: $+583\%$, weakest gap: $+300\%$.  The gold
    standard for threshold kinship detection.
  \item \textbf{Track~B (ambitious research):} The symbiogenesis family---four
    algorithm families exploring nonlinear, structure-first inheritance
    detection.  This paper documents the fourth and most ambitious family.
\end{itemize}


% ════════════════════════════════════════════════════════════════════
\section{Methods}
\label{sec:methods}

\subsection{Dataset}

We evaluate on the CEPH~1463 pedigree~\cite{ceph1463, zook2019giab,
1000genomes2015}, using Illumina Platinum Pedigree data (100\,bp paired-end,
$\sim$30$\times$ WGS) for six family members:

\begin{pedigree}
Generation 1 (Grandparents):
  Paternal: NA12889 (Pat.GF) + NA12890 (Pat.GM)
  Maternal: NA12891 (Mat.GF) + NA12892 (Mat.GM)

Generation 2 (Parents):
  NA12877 (Father) = child of NA12889 + NA12890
  NA12878 (Mother) = child of NA12891 + NA12892
\end{pedigree}

Ten pairwise comparisons: 4 parent--child, 2 spousal/in-law, 4 cross-lineage.
Each sample limited to 50{,}000 reads.  No alignment, variant calling, or
phasing is performed.

\subsection{Pipeline Overview}

The Spectral Ecology pipeline has six stages:

\begin{enumerate}
  \item \textbf{Module extraction.}  Raw reads $\to$ k-mer co-occurrence
    graphs $\to$ connected components $\to$ modules.  Parameters: $k = 15$,
    window $= 150$\,bp, min support $= 3$, min cohesion $= 0.25$.
    Produces 295--449 modules per sample.
  \item \textbf{Co-occurrence graph construction.}  For each read, identify
    all modules present.  For each module pair on the same read, increment
    co-occurrence count.  Normalize:
    $w(A,B) = \text{cooccur}(A,B) / \sqrt{\text{freq}(A) \cdot \text{freq}(B)}$.
    Threshold at $w > 0.01$.  Produces 611--1{,}410 edges.
  \item \textbf{Normalized graph Laplacian.}
    $L_{\text{sym}} = I - D^{-1/2} A D^{-1/2}$.
    Eigenvalues in $[0, 2]$; $\lambda_1 = 0$ for each connected component.
  \item \textbf{Jacobi eigendecomposition.}  Full dense eigendecomposition via
    Givens rotations.  Returns all eigenvalues (sorted ascending) and
    eigenvectors.  Sub-second for $n \leq 450$.
  \item \textbf{Ecosystem signature extraction.}  From the eigendecomposition:
    \begin{itemize}
      \item Heat kernel traces at 10 timescales ($t = 0.1$ to $100$)
      \item Per-node heat kernel signatures (HKS)
      \item Spectral ecological role assignment (core/satellite/bridge/scaffold/boundary)
      \item Eigenvalue density profile (Gaussian KDE, filtered for non-trivial eigenvalues)
      \item Eigenvector localization (inverse participation ratio)
      \item Diffusion distance distributions at 3 timescales
      \item Module feature embeddings (7-dimensional)
      \item Null model spectral residual (Chung-Lu baseline)
    \end{itemize}
  \item \textbf{Ecosystem comparison.}  Weighted combination of distances
    across all invariants.  Converted to similarity via
    $s = \exp(-d \cdot \alpha)$.
\end{enumerate}

\subsection{Module Co-occurrence Graph (v2)}

The critical design choice.  Version~1 used sequential adjacency (modules
appearing in order on the same read).  Version~2 uses statistical
co-occurrence:

\begin{definition}
Two modules $A$ and $B$ are connected if they co-occur on any read.  The edge
weight is the normalized co-occurrence:
\begin{equation}
  w(A, B) = \frac{|\{r : A \in r \wedge B \in r\}|}
                  {\sqrt{|\{r : A \in r\}| \cdot |\{r : B \in r\}|}}
  \label{eq:cooccurrence}
\end{equation}
Edges with $w < 0.01$ are removed.
\end{definition}

\ymarginnote{Co-occurrence is linkage disequilibrium without coordinates.}

This normalization is equivalent to cosine similarity in read-space.  It
approximates linkage disequilibrium without alignment coordinates: modules that
co-occur more often than expected by chance (given their individual frequencies)
receive high weight.

The impact is dramatic:

\begin{table}[H]
  \centering
  \tablestyle
  \caption{Graph density comparison: sequential adjacency (v1) vs.\ statistical
    co-occurrence (v2).}
  \label{tab:density}
  \begin{tabular}{lrrrr}
    \toprule
    \rowcolor{table-header}
    \textcolor{white}{\textbf{Sample}} &
    \textcolor{white}{\textbf{Modules}} &
    \textcolor{white}{\textbf{v1 Edges}} &
    \textcolor{white}{\textbf{v2 Edges}} &
    \textcolor{white}{\textbf{Density}} \\
    \midrule
    NA12889 (Pat.\ GF) & 373 & 2 & 951 & 1.4\% \\
    NA12890 (Pat.\ GM) & 295 & 4 & 611 & 1.4\% \\
    NA12891 (Mat.\ GF) & 449 & 5 & 1{,}410 & 1.4\% \\
    NA12892 (Mat.\ GM) & 402 & 5 & 1{,}067 & 1.3\% \\
    NA12877 (Father)    & 444 & 2 & 975 & 1.0\% \\
    NA12878 (Mother)    & 443 & 3 & 885 & 0.9\% \\
    \bottomrule
  \end{tabular}
\end{table}

Average edges per sample increased from 3.5 to 983---a \textbf{280$\times$
improvement}.

\subsection{Normalized Graph Laplacian}

Given the weighted adjacency matrix $A$ with degree matrix $D$
($D_{ii} = \sum_j A_{ij}$), the normalized Laplacian is:

\begin{equation}
  L_{\text{sym}} = I - D^{-1/2} A D^{-1/2}
  \label{eq:laplacian}
\end{equation}

Properties: symmetric positive semidefinite, eigenvalues in $[0, 2]$,
$\lambda_1 = 0$ for each connected component.  The spectral gap
$\lambda_2 - \lambda_1$ measures algebraic connectivity.

\subsection{Jacobi Eigendecomposition}

We implement the classical Jacobi rotation algorithm for symmetric
matrices~\cite{efron1979bootstrap}.  The algorithm iteratively applies Givens
rotations to zero off-diagonal elements, converging to the diagonal form
$\Lambda = V^T L V$ where $\Lambda$ contains eigenvalues and columns of $V$
are eigenvectors.

Key implementation detail: when diagonal elements are equal
($L_{pp} = L_{qq}$), the standard formula produces $\tau = 0$ and
$\text{sign}(0) = 0$ in JavaScript, making the rotation identity.  We
explicitly handle this case with $t = 1$ (rotation angle $\pi/4$).

For $n \leq 450$ nodes, full dense eigendecomposition completes in under one
second.

\subsection{Multi-Scale Heat Kernel}

The heat kernel $K_t = \exp(-tL)$ describes diffusion on the graph at
timescale $t$.  From the eigendecomposition $L = V \Lambda V^T$:

\begin{equation}
  K_t = V \exp(-t\Lambda) V^T
  \label{eq:heat-kernel}
\end{equation}

The per-node heat kernel signature (HKS) at node $j$ and time $t$:

\begin{equation}
  \text{HKS}(j, t) = \sum_{k=1}^{n} \exp(-t \lambda_k) \cdot \phi_k(j)^2
  \label{eq:hks}
\end{equation}

where $\phi_k(j)$ is the $j$-th component of eigenvector $k$.  Small $t$
captures local neighborhood structure; large $t$ captures global topology.

\ymarginnote{Heat kernel at multiple scales = multi-resolution fingerprint.}

The graph-level trace:

\begin{equation}
  \text{trace}(K_t) = \sum_{k=1}^{n} \exp(-t \lambda_k)
  \label{eq:trace}
\end{equation}

We evaluate at 10 timescales: $t \in \{0.1, 0.25, 0.5, 1.0, 2.0, 5.0, 10.0,
20.0, 50.0, 100.0\}$.

\subsection{Spectral Ecological Role Assignment}

Each module is classified into an ecological role based on its position in
spectral space:

\begin{definition}
The ecological roles are:
\begin{itemize}
  \item \textbf{Core}: high degree, central in spectral embedding.  Foundational
    modules around which the ecosystem organizes.
  \item \textbf{Satellite}: low degree, far from spectral center.  Peripheral
    modules loosely attached to the ecosystem.
  \item \textbf{Bridge}: high Fiedler value (2nd eigenvector component),
    central position.  Connects otherwise-disconnected clusters.
  \item \textbf{Scaffold}: moderate degree, moderate spectral position.
    Structural support modules.
  \item \textbf{Boundary}: high Fiedler value, far from center.  At the
    transition between spectral clusters.
\end{itemize}
\end{definition}

Roles are assigned from the first 6 non-trivial eigenvectors combined with
weighted graph degree.  Each node receives a confidence score based on the
margin between its highest and second-highest role score.

The \emph{role spectrum} is the distribution of roles across the ecosystem:
$(\text{frac}_{\text{core}}, \text{frac}_{\text{sat}}, \text{frac}_{\text{bridge}},
\text{frac}_{\text{scaffold}}, \text{frac}_{\text{boundary}})$.

\subsection{Additional Spectral Invariants}

\subsubsection{Eigenvalue Density}

Gaussian kernel density estimation over the non-trivial eigenvalues
(filtering out $\lambda < 0.01$ and $|\lambda - 1| < 0.01$).  32 bins over
$[0, 2]$.  Captures the shape of the connectivity spectrum.

\subsubsection{Eigenvector Localization}

The inverse participation ratio (IPR) for eigenvector $k$:
\begin{equation}
  \text{IPR}(k) = \sum_{j=1}^{n} \phi_k(j)^4
\end{equation}
IPR $= 1/n$ for delocalized (uniform) eigenvectors; IPR $= 1$ for fully
localized.  High IPR indicates modular structure.

\subsubsection{Diffusion Distance Distributions}

Pairwise diffusion distance between nodes $i$ and $j$ at time $t$:
\begin{equation}
  d_t(i,j)^2 = \sum_{k=1}^{n} \exp(-2t\lambda_k)(\phi_k(i) - \phi_k(j))^2
\end{equation}
We compute the mean and variance of all $\binom{n}{2}$ pairwise distances at
three timescales ($t = 0.5, 2.0, 10.0$).

\subsubsection{Node-Level HKS Distribution Comparison}

Rather than comparing heat traces (graph-level sums), we compare the
\emph{distribution} of per-node HKS values.  At each timescale, sort the
per-node HKS values for both samples, interpolate to a common grid, and compute
L2 distance between the sorted vectors.  This is equivalent to the 1-Wasserstein
distance between empirical node-signature distributions.

Before comparison, each node's HKS is normalized by the within-sample mean at
that timescale, so the distribution captures \emph{relative} node importance
(hub-ness vs.\ periphery) rather than absolute heat values.

\subsection{Ecosystem Comparison}

The composite distance between two ecosystem signatures is:

\begin{equation}
  d(A, B) = \sum_{i} w_i \cdot d_i(A, B)
  \label{eq:composite}
\end{equation}

where the components and weights (final configuration) are:

\begin{table}[H]
  \centering
  \tablestyle
  \caption{Comparison components and weights.}
  \label{tab:weights}
  \begin{tabular}{llc}
    \toprule
    \rowcolor{table-header}
    \textcolor{white}{\textbf{Component}} &
    \textcolor{white}{\textbf{Method}} &
    \textcolor{white}{\textbf{Weight}} \\
    \midrule
    Node HKS distribution & 1-Wasserstein (sorted L2) & 0.30 \\
    Heat trace & L2 of log-normalized traces & 0.10 \\
    Role spectrum & Jensen--Shannon divergence & 0.20 \\
    Eigenvalue density & L2 of KDE profiles & 0.15 \\
    Diffusion distribution & L2 of mean + variance & 0.10 \\
    Filtered spectral & L2 of non-trivial eigenvalues & 0.10 \\
    Trivial fraction & $|\text{frac}_A - \text{frac}_B|$ & 0.05 \\
    \bottomrule
  \end{tabular}
\end{table}

Similarity is computed as $s = \exp(-2d)$.


% ════════════════════════════════════════════════════════════════════
\section{Experimental Trajectory}
\label{sec:experiments}

We document six iterations, each changing one variable.

\subsection{v1: Sequential Adjacency (Baseline Failure)}

The initial implementation used \texttt{buildModuleGraph()} from
\texttt{@yalumba/modules}, which creates edges between modules that appear
\emph{sequentially} on the same read (module $A$'s motifs appear before module
$B$'s motifs, observed $\geq 2$ times).

\begin{limitation}
\textbf{Catastrophic graph collapse.}  With 295--449 modules per sample, only
2--5 edges survived the $\geq 2$ observation filter.  Graphs were $>$99\%
disconnected.  Spectral gap was exactly zero for all samples.  All eigenvalues
were either 0 (disconnected components) or 1 (isolated nodes).  The heat kernel
trace reduced to counting connected component sizes.  Role assignment was
trivial: 95\% satellite.
\end{limitation}

\begin{table}[H]
  \centering
  \tablestyle
  \caption{v1 results: sequential adjacency.}
  \label{tab:v1}
  \begin{tabular}{lr}
    \toprule
    \rowcolor{table-header}
    \textcolor{white}{\textbf{Metric}} &
    \textcolor{white}{\textbf{Value}} \\
    \midrule
    Mean separation & $-0.74\%$ \\
    Weakest gap & $-26.87\%$ \\
    Avg.\ edges per sample & 3.5 \\
    Spectral gap (all samples) & 0.0000 \\
    Dominant role & satellite (95\%) \\
    \bottomrule
  \end{tabular}
\end{table}

The spouse pair (Father~$\leftrightarrow$~Mother) ranked first at 0.90
similarity because they had the most similar module counts (444 vs.\ 443).
This is pure coverage confound.

\subsection{v2: Statistical Co-occurrence}

Replaced sequential adjacency with read-level co-occurrence
(Equation~\ref{eq:cooccurrence}).  Every module pair present on the same read
receives an edge, weighted by normalized frequency.

\begin{keyfinding}
\textbf{Separation flipped positive}: $+1.01\%$ (from $-0.74\%$).  Edges
increased 280$\times$.  The spectral machinery now had real structure to
analyze.
\end{keyfinding}

However, spectral gap remained zero (54 connected components, 44 isolated
nodes).  The graph was denser but still disconnected.

\subsection{Experiment 1: kNN Kernel Edges}

Added a support-similarity Gaussian kernel connecting each node to its 8
nearest neighbors by $\log$-support distance.  Guaranteed graph connectivity.

\begin{table}[H]
  \centering
  \tablestyle
  \caption{Experiment~1 results: kNN kernel edges.}
  \label{tab:exp1}
  \begin{tabular}{lrr}
    \toprule
    \rowcolor{table-header}
    \textcolor{white}{\textbf{Metric}} &
    \textcolor{white}{\textbf{Before}} &
    \textcolor{white}{\textbf{After}} \\
    \midrule
    Edges (avg.) & 983 & 3{,}624 \\
    Density (avg.) & 1.2\% & 5.3\% \\
    Spectral gap & 0.00 & 0.08--0.12 \\
    Separation & $+1.01\%$ & $-1.60\%$ \\
    \bottomrule
  \end{tabular}
\end{table}

\begin{limitation}
\textbf{Kernel edges amplify coverage confound.}  The support-similarity kernel
connects modules with similar read counts---a coverage proxy.  Samples with
similar sequencing depth get artificially similar graphs, worsening the coverage
confound.  Separation went \emph{backwards}.
\end{limitation}

The kernel was reverted.  Lesson: connectivity guarantees must use
coverage-orthogonal features.

\subsection{Experiment 2: Disable Broken Deconfounding}

The Chung-Lu null model produced spectral residuals with magnitudes of
65{,}000--240{,}000---numerically unstable z-scores caused by near-zero
variance in the null spectrum.  Even with $\tanh$ compression, the component
contributed a constant $\sim$0.10 to all pairs.

Removed the deconfounded component and redistributed its weight to heat trace.

\begin{keyfinding}
Separation improved to $+1.33\%$.  The broken component was consuming weight
budget without contributing discriminative signal.
\end{keyfinding}

\subsection{Experiment 3: Trivial Eigenvalue Filtering}

Observed that the eigenvalue spectrum is dominated by trivial values:

\begin{itemize}
  \item 12--13 eigenvalues near zero ($< 0.01$): disconnected components
  \item 72--117 eigenvalues near one ($|{\lambda - 1}| < 0.01$): isolated or
    loosely connected nodes
  \item 314--364 ``interesting'' eigenvalues in $(0.01, 0.99) \cup (1.01, 2.0)$
\end{itemize}

\ymarginnote{The interesting eigenvalues are only 70\% of the total.}

The near-one eigenvalues flatten the density estimate and dominate the spectral
distance.  Filtering them before density estimation and spectral comparison
isolates the structural signal.

\begin{keyfinding}
\textbf{Best separation achieved: $+1.89\%$.}  Eigenvalue filtering removed
the dominant noise from the spectral comparison, amplifying the structural
signal.
\end{keyfinding}

\subsection{Experiment 4: Node-Level HKS Distribution}

Replaced the graph-level heat trace (sum of node HKS) with node-level
distribution comparison.  At each timescale, sorted per-node HKS values are
compared via interpolated L2 distance---equivalent to 1-Wasserstein distance
between empirical distributions.

\begin{keyfinding}
Mat.GM~$\to$~Mother (parent--child) jumped to rank~2 for the first time.
The node-level distribution detects that grandmother and granddaughter have
similar internal hub/periphery structure, even though their ecosystems differ
in absolute size.
\end{keyfinding}

\subsection{Experiment 5: Mean-Normalized HKS}

Normalized each node's HKS by the within-sample mean at each timescale before
sorting and comparing.  This removes absolute scale, comparing only
\emph{relative} node importance.

Three of four parent--child pairs now rank consecutively at positions 3--5.

\begin{table}[H]
  \centering
  \tablestyle
  \caption{Experimental trajectory summary.}
  \label{tab:trajectory}
  \begin{tabular}{llcc}
    \toprule
    \rowcolor{table-header}
    \textcolor{white}{\textbf{Experiment}} &
    \textcolor{white}{\textbf{Change}} &
    \textcolor{white}{\textbf{Separation}} &
    \textcolor{white}{\textbf{Weakest Gap}} \\
    \midrule
    v1 & Sequential adjacency & $-0.74\%$ & $-26.87\%$ \\
    v2 & Co-occurrence edges & $+1.01\%$ & $-17.72\%$ \\
    Exp.\ 1 & kNN kernel & $-1.60\%$ & $-13.13\%$ \\
    Exp.\ 2 & Drop deconfounding & $+1.33\%$ & $-26.73\%$ \\
    Exp.\ 3 & Filter eigenvalues & $\mathbf{+1.89\%}$ & $-26.75\%$ \\
    Exp.\ 4 & Node HKS distribution & $+1.35\%$ & $-21.66\%$ \\
    Exp.\ 5 & Mean-normalized HKS & $+1.55\%$ & $-22.53\%$ \\
    \bottomrule
  \end{tabular}
\end{table}


% ════════════════════════════════════════════════════════════════════
\section{Results}
\label{sec:results}

\subsection{Final Score Table}

\begin{table}[H]
  \centering
  \tablestyle
  \caption{Complete pairwise scores (Experiment~5 configuration), ranked by
    similarity.  Parent--child pairs marked with $\blacktriangleright$.}
  \label{tab:scores}
  \begin{tabular}{clcl}
    \toprule
    \rowcolor{table-header}
    \textcolor{white}{\textbf{Rank}} &
    \textcolor{white}{\textbf{Pair}} &
    \textcolor{white}{\textbf{Similarity}} &
    \textcolor{white}{\textbf{Type}} \\
    \midrule
    1 & Father $\leftrightarrow$ Mother
      & 0.803 & Spouses \\
    2 & Pat.GF $\leftrightarrow$ Mother
      & 0.749 & In-law \\
    3 & $\blacktriangleright$ Mat.GM $\to$ Mother
      & 0.730 & Parent--child \\
    4 & $\blacktriangleright$ Pat.GM $\to$ Father
      & 0.725 & Parent--child \\
    5 & $\blacktriangleright$ Pat.GF $\to$ Father
      & 0.715 & Parent--child \\
    6 & Pat.GM $\leftrightarrow$ Mat.GM
      & 0.692 & Unrelated \\
    7 & Pat.GF $\leftrightarrow$ Mat.GM
      & 0.687 & Unrelated \\
    8 & $\blacktriangleright$ Mat.GF $\to$ Mother
      & 0.578 & Parent--child \\
    9 & Mat.GF $\leftrightarrow$ Father
      & 0.566 & In-law \\
    10 & Pat.GF $\leftrightarrow$ Mat.GF
      & 0.530 & Unrelated \\
    \bottomrule
  \end{tabular}
\end{table}

\subsection{Separation and Gap}

\begin{table}[H]
  \centering
  \tablestyle
  \caption{Summary metrics.}
  \label{tab:summary}
  \begin{tabular}{lr}
    \toprule
    \rowcolor{table-header}
    \textcolor{white}{\textbf{Metric}} &
    \textcolor{white}{\textbf{Value}} \\
    \midrule
    Parent--child average & 0.687 \\
    Unrelated average & 0.671 \\
    Mean separation & $+1.55\%$ \\
    Weakest gap & $-22.53\%$ \\
    Detects all & NO \\
    Prep time & 379\,s \\
    Compare time & $<0.01$\,s \\
    Heap delta & $+$1{,}853\,MB \\
    \bottomrule
  \end{tabular}
\end{table}

\subsection{Per-Sample Ecosystem Profiles}

\begin{table}[H]
  \centering
  \tablestyle
  \caption{Ecosystem signatures for all six samples.}
  \label{tab:ecosystems}
  \begin{tabular}{llrrcccc}
    \toprule
    \rowcolor{table-header}
    \textcolor{white}{\textbf{Sample}} &
    \textcolor{white}{\textbf{Role}} &
    \textcolor{white}{\textbf{Mods}} &
    \textcolor{white}{\textbf{Edges}} &
    \textcolor{white}{\textbf{Core}} &
    \textcolor{white}{\textbf{Sat}} &
    \textcolor{white}{\textbf{Bridge}} &
    \textcolor{white}{\textbf{Entropy}} \\
    \midrule
    NA12889 & Pat.\ GF & 373 & 951  & 230 & 141 & 0   & 5.74 \\
    NA12890 & Pat.\ GM & 295 & 611  & 198 & 93  & 1   & 5.52 \\
    NA12891 & Mat.\ GF & 449 & 1410 & 8   & 85  & 356 & 5.95 \\
    NA12892 & Mat.\ GM & 402 & 1067 & 296 & 104 & 0   & 5.81 \\
    NA12877 & Father   & 444 & 975  & 325 & 118 & 0   & 5.91 \\
    NA12878 & Mother   & 443 & 885  & 311 & 130 & 0   & 5.92 \\
    \bottomrule
  \end{tabular}
\end{table}

\begin{keyfinding}
\textbf{NA12891 (Mat.~GF) is a structural outlier.}  It has 356 bridges and
only 8 cores---a completely inverted role distribution compared to every other
sample (which have 198--325 cores and 0 bridges).  This produces a role
spectrum Jensen--Shannon divergence of $\sim$0.71 for all pairs involving
NA12891, dragging Mat.GF~$\to$~Mother to rank~8 and Mat.GF~$\leftrightarrow$~Father
to rank~9.
\end{keyfinding}

\subsection{Eigenvalue Spectrum Analysis}

The eigenvalue spectrum of each sample's co-occurrence graph reveals three
distinct populations:

\begin{itemize}
  \item \textbf{Near-zero} ($\lambda < 0.01$): 12--13 eigenvalues per sample,
    corresponding to disconnected components in the co-occurrence graph.
  \item \textbf{Near-one} ($|\lambda - 1| < 0.01$): 72--117 eigenvalues,
    corresponding to isolated or loosely connected modules.
  \item \textbf{Structural} ($\lambda \in (0.01, 0.99) \cup (1.01, 2.0)$):
    314--364 eigenvalues carrying the actual connectivity information.
\end{itemize}

Father and Mother have the most similar eigenvalue breakdowns: 13/101/330 vs.\
12/117/314 near-zero/near-one/structural.  Mat.~GF has fewer near-one
eigenvalues (72) and more structural (364), reflecting its denser, more
interconnected graph.

\subsection{Comparison to All Symbiogenesis Algorithms}

\begin{table}[H]
  \centering
  \tablestyle
  \caption{All symbiogenesis algorithm families on CEPH~1463.}
  \label{tab:all-algorithms}
  \begin{tabular}{lcccc}
    \toprule
    \rowcolor{table-header}
    \textcolor{white}{\textbf{Algorithm}} &
    \textcolor{white}{\textbf{Separation}} &
    \textcolor{white}{\textbf{Weakest Gap}} &
    \textcolor{white}{\textbf{Time}} &
    \textcolor{white}{\textbf{Status}} \\
    \midrule
    Rare-Run P90
      & $+$583\% & $+$300\% & 72.5\,s & \textbf{Gold} \\
    Module Persistence v3
      & $+$3.36\% & $+$0.13\% & 1{,}085\,s & \textbf{PASSES} \\
    Coalition Transfer v3
      & $+$1.97\% & $+$0.06\% & 272\,s & \textbf{PASSES} \\
    \textbf{Spectral Ecology v2}
      & $+$1.55\% & $-$22.53\% & 379\,s & Fails \\
    Module Stability v4
      & $+$4.90\% & $-$1.04\% & 405\,s & Fails \\
    \bottomrule
  \end{tabular}
\end{table}

Spectral Ecology achieves competitive separation ($+1.55\%$) but the worst
weakest gap ($-22.53\%$) among the symbiogenesis family.  The gap reflects the
two outlier pairs; without them, the remaining 8 pairs would separate cleanly.


% ════════════════════════════════════════════════════════════════════
\section{Failure Analysis}
\label{sec:failure}

\subsection{Failure Mode 1: The Spouse Confound}

Father (NA12877) and Mother (NA12878) consistently rank as the most similar
pair across all spectral experiments.  Their ecosystem signatures are nearly
identical:

\begin{itemize}
  \item Module count: 444 vs.\ 443
  \item Edge count: 975 vs.\ 885
  \item Role distribution: 325C/118S vs.\ 311C/130S
  \item Spectral entropy: 5.913 vs.\ 5.915
  \item Localization: 0.203 vs.\ 0.208
\end{itemize}

\ymarginnote{The spouses are spectrally indistinguishable from parent--child
pairs.}

Every graph-level metric and every distribution-level metric sees them as
twins.  This is because:

\begin{enumerate}
  \item They have nearly identical module counts (similar genomic complexity at
    this read depth).
  \item Their co-occurrence graphs have similar density and degree distributions.
  \item The spectral decomposition of two graphs with similar size and density
    produces similar eigenvalue distributions.
  \item Role assignment depends on the spectral embedding, which is similar.
\end{enumerate}

\begin{limitation}
\textbf{Purely within-sample spectral methods cannot distinguish ``similar
because related'' from ``similar because similar sequencing.''} Both produce
similar graph topology.  The spouse pair demonstrates that graph-level and
distribution-level spectral invariants are necessary but not sufficient for
relatedness detection.
\end{limitation}

The fix requires \emph{cross-sample} signal: detecting that specific
substructures (not just aggregate statistics) are shared between individuals.

\subsection{Failure Mode 2: The NA12891 Outlier}

NA12891 (Mat.~GF) has a unique role distribution: 8 cores, 85 satellites, 356
bridges.  Every other sample has 0 bridges.

This anomaly produces role spectrum JSD $\sim$0.71 for all NA12891 pairs,
compared to $< 0.03$ for pairs not involving NA12891.  Since role spectrum has
20\% weight, this alone contributes $0.20 \times 0.71 = 0.14$ to the distance
for NA12891 pairs.

\ymarginnote{One outlier sample can dominate the role spectrum signal.}

The bridge classification arises because NA12891 has the densest graph (1{,}410
edges, 1.4\% density) with many modules at high Fiedler values and central
spectral positions.  This may reflect:

\begin{itemize}
  \item A genuine biological difference (different haplotype block structure)
  \item A sequencing library artifact (different insert size distribution)
  \item A sensitivity of the role classifier to graph density
\end{itemize}

Robust role assignment would need to normalize for graph density before
classifying.


% ════════════════════════════════════════════════════════════════════
\section{Discussion}
\label{sec:discussion}

\subsection{What Spectral Ecology Achieved}

The spectral approach proved three things:

\begin{enumerate}
  \item \textbf{Dense co-occurrence graphs are tractable and informative.}  The
    switch from sequential adjacency (2--5 edges) to statistical co-occurrence
    (600--1{,}400 edges) transformed useless disconnected graphs into spectrally
    rich structures with diverse eigenvalue distributions and meaningful role
    assignments.
  \item \textbf{Identity-free comparison works in the right direction.}  With
    $+1.55\%$ separation and three of four parent--child pairs at ranks 3--5,
    the spectral method detects that related individuals produce more similar
    ecosystem organization than unrelated individuals---without any token
    matching.
  \item \textbf{Node-level distribution comparison outperforms graph-level
    summaries.}  Comparing the distribution of per-node heat kernel signatures
    (1-Wasserstein) captures internal hub/periphery structure that the trace
    (just the sum) misses.
\end{enumerate}

\subsection{The Within-Sample Ceiling}

All experiments converge to $\sim$1.5--1.9\% separation, regardless of which
invariants are used or how they are weighted.  This suggests a fundamental
ceiling for within-sample spectral methods.

The ceiling exists because within-sample spectral features characterize the
\emph{shape} of the ecosystem---its connectivity pattern, degree distribution,
cluster structure.  But shape similarity correlates with size similarity
(similar $n$, similar density), which correlates with coverage similarity.
The coverage confound enters through the graph construction, not the spectral
analysis.

\begin{limitation}
\textbf{The spectral approach correctly measures what it measures.}  The
problem is that what it measures---ecosystem shape---is correlated with
coverage.  Breaking through requires either (a)~explicit coverage
deconfounding that works (the Chung-Lu null model failed) or (b)~cross-sample
structural signal that is coverage-orthogonal.
\end{limitation}

\subsection{The Path Forward}

The ChatGPT consultation identified the core issue as ``information starvation
in the graph construction layer'' and prescribed a multi-layer ecosystem graph.
Our experiments confirmed this diagnosis.  Three directions are most promising:

\subsubsection{Multi-Layer Ecosystem Graph}

Build four coupled graphs:
\begin{enumerate}
  \item \textbf{Motif compatibility}: normalized mutual information between
    k-mer co-occurrence patterns.
  \item \textbf{Module similarity}: cosine/JSD between motif frequency
    distributions---soft matching without identity.
  \item \textbf{Coalition affinity}: read-level co-occurrence (implemented
    in v2).
  \item \textbf{Perturbation response}: correlation of bootstrap survival
    probability, with coverage regressed out.
\end{enumerate}

The weighted sum Laplacian $L = \alpha L_1 + \beta L_2 + \gamma L_3 + \delta L_4$
would capture multiple ecological scales simultaneously.

\subsubsection{Cross-Sample Module Embedding Similarity}

Embed modules from different samples into a shared feature space (Section~V of
the ChatGPT response).  Two modules with zero k-mer overlap can be
``ecologically equivalent'' if they occupy similar positions in the embedding.
Compare samples via optimal transport between their embedding distributions.

\subsubsection{Topological Methods}

Build simplicial complexes from higher-order co-occurrence (triples, quadruples
of modules on the same read) and compute persistent homology.  Topological
invariants (Betti numbers, persistence diagrams) are robust to small
perturbations and could provide coverage-orthogonal features.

\subsection{Connection to the Symbiogenesis Thesis}

The symbiogenesis thesis predicts that inheritance transmits organizational
constraints, not specific tokens.  Spectral Ecology is the first algorithm to
test this prediction directly: it measures organizational similarity without
token matching.

The result is encouraging.  Parent--child pairs \emph{do} have more similar
ecosystem organization than most unrelated pairs ($+1.55\%$ separation).  The
organizational signal is real---it's just weaker than the coverage confound.

\begin{pullquote}
The spectral approach finds the signal. The challenge is that the noise
(coverage confound) is also structural.
\end{pullquote}

This suggests that the thesis is correct but the implementation needs richer
structure.  The organization that is inherited is not just graph shape---it
includes which specific \emph{types} of modules interact, which coalitions
persist under perturbation, and how the ecosystem recovers from disruption.
Capturing these requires the multi-layer, dynamical, cross-sample methods
outlined above.


% ════════════════════════════════════════════════════════════════════
\section{Implementation}
\label{sec:implementation}

The \texttt{@yalumba/ecology} package contains eight source modules:

\begin{table}[H]
  \centering
  \tablestyle
  \caption{Package structure of \texttt{@yalumba/ecology}.}
  \label{tab:package}
  \begin{tabular}{ll}
    \toprule
    \rowcolor{table-header}
    \textcolor{white}{\textbf{Module}} &
    \textcolor{white}{\textbf{Responsibility}} \\
    \midrule
    \texttt{motif-graph.ts} & Co-occurrence graph construction \\
    \texttt{laplacian.ts} & Normalized Laplacian + Jacobi eigendecomposition \\
    \texttt{heat-kernel.ts} & Multi-scale heat kernel signatures \\
    \texttt{role-assignment.ts} & Spectral ecological role classification \\
    \texttt{spectral-invariants.ts} & Eigenvalue density, localization, curvature, diffusion \\
    \texttt{null-model.ts} & Chung-Lu null model + spectral residual \\
    \texttt{module-embedding.ts} & Module feature vectors + cross-sample similarity \\
    \texttt{compare.ts} & Composite ecosystem comparison \\
    \bottomrule
  \end{tabular}
\end{table}

All code is TypeScript on the Bun runtime.  No external numerical libraries.
The Jacobi eigendecomposition, heat kernel computation, and Gaussian KDE are
all implemented from scratch.

Dependencies: \texttt{@yalumba/modules} (module extraction, co-occurrence),
\texttt{@yalumba/math} (statistics, information theory).

Ten tests verify the full pipeline: Laplacian eigenvalue properties, heat
kernel monotonicity, role assignment completeness, signature construction,
and comparison transitivity.


% ════════════════════════════════════════════════════════════════════
\section{Limitations}
\label{sec:limitations}

\begin{limitation}
\textbf{Negative weakest gap.}  The $-22.53\%$ gap means the method cannot
classify any pair correctly at any threshold.  It fails as a detector.
\end{limitation}

\begin{limitation}
\textbf{Coverage confound persists.}  The spouse pair ranks first because
similar sequencing depth produces similar ecosystem structure.  Within-sample
spectral methods cannot distinguish this from genuine relatedness.
\end{limitation}

\begin{limitation}
\textbf{NA12891 outlier.}  One sample's anomalous role distribution dominates
two of ten pairwise comparisons.  The method is not robust to per-sample data
quality variation.
\end{limitation}

\begin{limitation}
\textbf{Null model failure.}  The Chung-Lu random graph model produces
numerically unstable spectral residuals (z-scores in the 100{,}000s).  A
proper configuration model or degree-preserving rewiring is needed.
\end{limitation}

\begin{limitation}
\textbf{No cross-sample signal.}  All features are computed per-sample then
compared.  No information about whether specific structural motifs are
\emph{shared} across samples---only whether aggregate statistics are similar.
\end{limitation}

\begin{limitation}
\textbf{Single dataset.}  Six European-ancestry individuals, 10 pairs.
Statistical power is minimal.
\end{limitation}

\begin{limitation}
\textbf{High memory.}  The Jacobi eigendecomposition stores dense $n \times n$
matrices, consuming 1{,}853\,MB heap.  Sparse methods would reduce this for
larger graphs.
\end{limitation}

\begin{limitation}
\textbf{Graph still sparse.}  At 1.0--1.4\% density and spectral gap $= 0$,
the co-occurrence graph remains mostly disconnected.  The ``interesting''
eigenvalues (those between 0.01 and 0.99) are only 70\% of the total.
\end{limitation}


% ════════════════════════════════════════════════════════════════════
\section{Conclusion}
\label{sec:conclusion}

Spectral Ecology is the first symbiogenesis algorithm that measures
\emph{organizational similarity without token identity}.  The switch from
sequential adjacency to statistical co-occurrence transformed unusable
disconnected graphs into spectrally rich structures, and five iterative
experiments pushed separation from $-0.74\%$ to $+1.55\%$---confirming that
related individuals produce more similar ecosystem organization than unrelated
individuals.

\ymarginnote{The spectral signal is real. The challenge is extracting it from
the coverage noise.}

The method's $-22.53\%$ weakest gap confirms that within-sample spectral
invariants alone are insufficient for threshold-based relatedness detection.
Two failure modes---the spouse coverage confound and the NA12891 outlier---
account for the entire gap.  Three of four parent--child pairs rank
consecutively at positions 3--5, demonstrating that the spectral approach
works qualitatively even where it fails quantitatively.

The path forward is clear: cross-sample structural signal, multi-layer
ecosystem graphs, and coverage-orthogonal features.  The spectral ecology
package provides the computational foundation---eigendecomposition, heat
kernels, role assignment, eigenvalue density, localization, diffusion
distances---on which these richer methods can build.

\begin{pullquote}
Symbiogenesis predicts that structure persists even when sequence changes.
The spectral approach finds structural similarity between related individuals.
The challenge is that coverage similarity produces the same signal.  Breaking
through requires methods that detect \emph{shared} structure, not just
\emph{similar} structure.
\end{pullquote}


% ════════════════════════════════════════════════════════════════════
\bibliographystyle{plain}
\bibliography{../references}

\end{document}
